% Independent completion of the unrestricted collar question, by /root/algebra_c.
% AI-produced by GPT-6 Astra (OpenAI), ChatGPT Pro mode, 19 November 2026.
\subsection{Arbitrary collars: there is no pointwise metric monotonicity}
\label{C:A2e-section}

The official scan~67/021 asks for the relation between the canonical metrics after removing an open ``collar''; it does not require that the new boundary be an equidistant curve. Proposition~\ref{A:collar} gives the exact nonlinear comparison and the strict increase for equidistant collars. The following result settles the unrestricted smooth, and even real-analytic, case.

\begin{theorem}[Proved: arbitrarily small analytic collars give incomparable metrics]
\label{C:A2e-nonmonotonicity}
Let $(X,h)$ be any compact connected hyperbolic surface with nonempty geodesic boundary and negative Euler characteristic. There are arbitrarily thin open collars of its boundary, with real-analytic inner boundaries, such that the remaining surface $X_t$ is diffeomorphic to $X$ but its canonical curvature $-1$ metric $h_t'$ with geodesic boundary is neither pointwise greater than nor pointwise less than $h|_{X_t}$. More precisely, their conformal ratio satisfies
\[
 h_t'=e^{2u_t}h|_{X_t},
 \qquad \min_{X_t}u_t<0<\max_{X_t}u_t.
\]
This includes every compact hyperbolic pair of pants and every compact one-holed torus. The new boundaries can converge to the original ones in $C^m$ for every fixed $m$.
\end{theorem}

\begin{proof}
Fix one boundary component $B$, of length $L$, and a disjoint Fermi collar $0\le r<R$ around it, with arclength coordinate $s\in\mathbb R/L\mathbb Z$. Here
\[
 h=dr^2+\cosh^2r\,ds^2,
 \qquad n|_B=-\partial_r.
\]
Choose $0<\varepsilon<1$, and, for a positive integer $N$ to be chosen, put $k=2\pi N/L$ and
\[
 f(s)=1+\varepsilon\cos(ks)\quad\text{on }B,
 \qquad f=1\quad\text{on every other boundary component}.
\]
Remove the relatively open collars $0\le r<t f(s)$, using Fermi coordinates at every boundary component. Since $f\ge1-\varepsilon>0$, these are collars of strictly positive width. For each fixed $N$ and all sufficiently small $t>0$, their closures are disjoint, their inner boundaries are analytic, and their complement $X_t$ is diffeomorphic to $X$.

We first derive the linear response rigorously, including the dependence on the moving domain. Choose smooth diffeomorphisms $T_t:X\to X_t$ which take $(0,s)$ to $(tf(s),s)$ and are the identity outside the fixed collars; for example, near a boundary take $r\mapsto r+t\chi(r)f(s)$ with $\chi=1$ near $0$ and $\chi=0$ farther inside. Put $\widehat h_t=T_t^*h$ and $U_t=u_t\circ T_t$. Then $\widehat h_t\to h$ smoothly, with difference $O(t)$ in every fixed differentiability norm. Proposition~\ref{A:collar} becomes
\begin{equation}
 \Delta_{\widehat h_t}U_t=e^{2U_t}-1,
 \qquad \partial_{n_t}U_t=-\kappa_t,
 \label{C:eq-collar-pulled}
\end{equation}
where $\kappa_t$ is the $h$-geodesic curvature of $r=tf(s)$ toward the remaining surface, pulled back to the original boundary. Direct differentiation in the displayed Fermi metric gives
\begin{equation}
 \kappa_t=t(f''-f)+O(t^2)
 \label{C:eq-collar-curvature}
\end{equation}
in every fixed differentiability norm. Indeed the radial component of the acceleration of $s\mapsto(tf(s),s)$ is $t f''-\cosh(tf)\sinh(tf)=t(f''-f)+O(t^2)$, while its unit tangent and inward unit normal have the limiting values $\partial_s$ and $\partial_r$; the tangential acceleration contributes only order $t^2$ to the curvature.

For clarity, the linear boundary problem used below is
\begin{equation}
 (\Delta_h-2)w=0,\qquad \partial_nw=b.
 \label{C:eq-collar-linear}
\end{equation}
It has a unique smooth solution for every smooth boundary function $b$. Its weak formulation is
\[
 \mathcal E(w,v):=\int_X
       (\langle\nabla w,\nabla v\rangle_h+2wv)\,dA_h
   =\int_{\partial X}bv\,ds_h.
\]
The trace theorem makes the right side bounded on $H^1(X)$, and $\mathcal E$ is coercive there, so the Hilbert-space representation theorem gives existence and uniqueness. Elliptic boundary regularity gives smoothness. The same linear problem, including its existence and estimate, appears in Rupflin (2021), Lemma~2.9(ii), equation~(2.22), pp.~3126--3127, specialized to geodesic boundary~\cite{C_Rupflin2021}. Unlike a pure Neumann Laplace equation, the term $-2w$ removes the constant-function kernel, so no compatibility condition is needed.

We show that the solution of~\eqref{C:eq-collar-pulled} obeys
\begin{equation}
 U_t/t\longrightarrow w_f\quad\text{in }H^1(X),
 \qquad \partial_nw_f=f-f''.
 \label{C:eq-collar-response}
\end{equation}
This argument does not presume a shape derivative theorem. Let $w_0$ solve~\eqref{C:eq-collar-linear} with $b=1$ on the entire boundary. The maximum principle, or the negative-part test in its weak equation, gives $w_0>0$. Set $P=w_0+1$. For a sufficiently large fixed constant $A$ and sufficiently small $t>0$, the functions $\pm AtP$ are upper and lower barriers for~\eqref{C:eq-collar-pulled}. To verify this explicitly, use
\[
 (\Delta_{\widehat h_t}-2)P=-2+O(t),
 \qquad \partial_{n_t}P=1+O(t),
 \qquad \kappa_t=O(t).
\]
Thus $\Delta_{\widehat h_t}(AtP)-(e^{2AtP}-1)<0$, its outward derivative is at least $-\kappa_t$, and the reversed inequalities hold for $-AtP$. Comparison follows by multiplying the equation for the positive part of the difference by that positive part and integrating; the interior zeroth-order coefficient is strictly positive and the boundary term has the opposite sign. Consequently $\|U_t\|_\infty\le Ct$.

Let $\mathcal E_t$ be the preceding energy form with metric $\widehat h_t$. The weak equation is
\[
 \mathcal E_t(U_t,v)
 =-\int_{\partial X}\kappa_t v\,ds_{\widehat h_t}
  -\int_X(e^{2U_t}-1-2U_t)v\,dA_{\widehat h_t}.
\]
Since $\|U_t\|_\infty=O(t)$, the second integrand is uniformly $O(t^2)|v|$. Uniform coercivity, the trace theorem, and the choice $v=U_t$ give $\|U_t\|_{H^1}=O(t)$. Divide the equation by $t$, subtract the weak equation with datum $f-f''$, and use~\eqref{C:eq-collar-curvature}. The differences of the metric coefficients and measures are $O(t)$, and the exponential remainder divided by $t$ is $O(t)$. We obtain
\[
 |\mathcal E(U_t/t-w_f,v)|\le Ct\|v\|_{H^1}.
\]
Taking $v=U_t/t-w_f$ proves~\eqref{C:eq-collar-response}, in fact with error $O(t)$. In particular, its boundary traces converge in $L^2$.

We now choose $N$ so that this first-order response is negative somewhere. Let $w_k$ solve~\eqref{C:eq-collar-linear} with datum $(1+k^2)\cos(ks)$ on $B$ and zero on all other boundaries. Then
\[
 w_f=w_0+\varepsilon w_k.
\]
Let $\phi_B$ solve the same linear problem with datum $1$ on $B$ and zero elsewhere. Reciprocity of $\mathcal E$ gives the exact identities
\begin{align}
 (1+k^2)\int_Bw_k\cos(ks)\,ds
   &=\mathcal E(w_k,w_k),\label{C:eq-collar-fourier-energy}\\
 \int_Bw_k\,ds
   &=(1+k^2)\int_B\phi_B\cos(ks)\,ds.
   \label{C:eq-collar-mean}
\end{align}
Set $K_3=\int_B|\partial_s^3\phi_B|\,ds<\infty$. Three integrations by parts on the closed boundary circle show that, for $k\ge1$,
\begin{equation}
 \left|\int_Bw_k\,ds\right|
 \le (1+k^2)K_3/k^3\le2K_3/k.
 \label{C:eq-collar-mean-bound}
\end{equation}

For the opposite estimate on its oscillating coefficient, use the explicit $H^1$ test function
\[
 v_k(r,s)=(1-r/R)e^{-kr}\cos(ks)\quad(0\le r\le R),
 \qquad v_k=0\quad\text{outside this collar}.
\]
It is continuous at $r=R$, so its zero extension lies in $H^1(X)$. Put $\chi(r)=1-r/R$. Direct integration of the angular variables gives
\[
\begin{aligned}
 \mathcal E(v_k,v_k)
 &=\frac L2\int_0^R e^{-2kr}
       \bigl[\cosh r\,(\chi'-k\chi)^2\\
 &\hspace{6em}+(k^2\operatorname{sech}r+2\cosh r)\chi^2\bigr]\,dr.
\end{aligned}
\]
For $k\ge\max(1,R^{-1})$, the bounds $(a+b)^2\le2a^2+2b^2$, $0\le1-r/R\le1$, and $\int_0^R e^{-2kr}dr\le(2k)^{-1}$ yield
\[
 \mathcal E(v_k,v_k)\le C_R k,
 \qquad C_R=\frac{L(6\cosh R+1)}4.
\]
But $v_k|_B=\cos(ks)$, and hence
\[
 \mathcal E(w_k,v_k)=(1+k^2)L/2.
\]
Cauchy--Schwarz for the positive inner product $\mathcal E$, followed by~\eqref{C:eq-collar-fourier-energy}, proves
\begin{equation}
 \int_Bw_k\cos(ks)\,ds
 \ge\frac{(1+k^2)L^2}{4C_Rk}
 \ge c_Rk,
 \qquad c_R=\frac{L}{6\cosh R+1}>0.
 \label{C:eq-collar-amplitude}
\end{equation}
No asymptotic expansion of a boundary pseudodifferential operator is used.

Let $M_0=\int_Bw_0\,ds>0$. Since $w_0>0$, its cosine coefficient has absolute value at most $M_0$. Equations~\eqref{C:eq-collar-mean-bound} and~\eqref{C:eq-collar-amplitude} imply
\[
 \int_B w_f(1-\cos(ks))\,ds
 \le2M_0+2\varepsilon K_3/k-\varepsilon c_Rk.
\]
Choose the integer $N$ so that $k=2\pi N/L$ satisfies
\[
 k>\max\left(1,R^{-1},\frac{4M_0}{\varepsilon c_R},
                        \sqrt{\frac{4K_3}{c_R}}\right).
\]
The displayed integral is then strictly negative. Its weight $1-\cos(ks)$ is nonnegative, so $w_f$ is negative at some point of $B$.

By~\eqref{C:eq-collar-response} and trace convergence,
\[
 \int_B U_t(1-\cos(ks))\,ds<0
\]
for all sufficiently small positive $t$. Therefore $U_t<0$ at some boundary point. Smoothness gives $u_t<0$ also on a nonempty relatively open region of $X_t$, including interior points, proving $h_t'<h$ there.

There must also be a region with the opposite inequality. Gauss--Bonnet and the fact that both canonical metrics have geodesic boundary give
\[
 \operatorname{Area}_{h_t'}(X_t)
  =-2\pi\chi(X_t)=-2\pi\chi(X)
  =\operatorname{Area}_h(X)
  >\operatorname{Area}_h(X_t).
\]
Hence $\int_{X_t}(e^{2u_t}-1)\,dA_h>0$, so $u_t>0$ somewhere. This proves incomparability. Finally $N$ is fixed before $t\downarrow0$; all derivatives of $tf(s)$ tend to zero. Thus the collars are arbitrarily thin and their analytic boundaries converge in every $C^m$ norm as asserted.
\end{proof}

\paragraph{Exact scope of the answer to 67/021.}
The nonlinear Neumann equation in Proposition~\ref{A:collar} characterizes the comparison for any smooth collar. Equidistant inward cuts strictly increase the metric. General analytic inward cuts can make the metrics incomparable, by Theorem~\ref{C:A2e-nonmonotonicity}. The manuscript itself asks a question and does not assert a false monotonicity theorem. Our counterexample refutes that potential generalization, while retaining the stated hypotheses of the question.
